\documentclass[10pt,a4paper]{article}
\usepackage[margin=18mm,headheight=14pt]{geometry}
\usepackage{fontspec}
\setmainfont{lmroman10-regular.otf}[BoldFont=lmroman10-bold.otf,ItalicFont=lmroman10-italic.otf,BoldItalicFont=lmroman10-bolditalic.otf]
\setsansfont{lmsans10-regular.otf}[BoldFont=lmsans10-bold.otf]
\setmonofont{lmmonolt10-regular.otf}[BoldFont=lmmonolt10-bold.otf,ItalicFont=lmmonolt10-oblique.otf]
\usepackage[ngerman]{babel}
\usepackage{amsmath,amssymb,booktabs,tabularx,array,fancyhdr,listings,xcolor}
\usepackage[hidelinks]{hyperref}
\pagestyle{fancy}\fancyhf{}
\fancyhead[L]{\sffamily AI-Hardware · Tag 4}
\fancyhead[R]{\sffamily Abschlussarbeit · 100 Punkte}
\fancyfoot[L]{\small Name: \rule{40mm}{0.3pt}}
\fancyfoot[R]{\thepage\ / 6}
\setlength{\parindent}{0pt}\setlength{\parskip}{5pt}
\setlength{\tabcolsep}{7pt}\renewcommand{\arraystretch}{1.18}
\lstset{language=Python,basicstyle=\ttfamily\small,
 keywordstyle=\bfseries,commentstyle=\itshape\color{black},
 stringstyle=\ttfamily\color{black},showstringspaces=false,
 frame=single,rulecolor=\color{black},framesep=5pt,
 columns=fullflexible,keepspaces=true,breaklines=true,tabsize=4,
 aboveskip=8pt,belowskip=8pt}
\newcommand{\titlepagepart}[2]{\vspace*{-4mm}{\sffamily\small #1}\par{\sffamily\LARGE\bfseries #2}\par\vspace{2mm}}
\newcommand{\block}[1]{\par\vspace{2mm}{\sffamily\large\bfseries #1}\par}
\newcommand{\code}[1]{\texttt{#1}}

\newcommand{\antwort}[1]{\par\vspace{#1}\hrule\vspace{2mm}}
\newcommand{\aufgabe}[2]{\block{#1 \hfill #2 Punkte}}
\begin{document}
\titlepagepart{01 / VERSTEHEN UND UMRECHNEN}{Abschlussarbeit · Tag 4}
\textbf{Name:} \rule{65mm}{0.3pt}\hfill\textbf{Datum:} \rule{30mm}{0.3pt}

\textbf{Umfang:} 100 Punkte, etwa 120 Minuten als Orientierung, ohne Zeitdruck.
Seiten 1--4 zuerst ohne Python bearbeiten; danach darfst du mit dem Cheatsheet
nacharbeiten. Ergänzungen kenntlich machen. Auf Seiten 5--6 sind Python,
Terminal und Cheatsheet erlaubt. Rechenwege zählen auch bei falschen Endergebnissen.
Leerzeichen in Bitmustern dienen nur der Lesbarkeit. Alle signed Zahlen sind
Zweierkomplementzahlen. Es gibt keine versteckten Lösungen in diesem Blatt.

\aufgabe{1. Das Warum erklären}{8}
\textbf{a) (4 P)} Ein Vier-Bit-Register enthält \code{1110}. Warum reicht das
Bitmuster allein nicht aus, um seine Bedeutung zu bestimmen? Nenne die Werte
unsigned und signed und erkläre, welche Vereinbarung den Unterschied macht.
\antwort{22mm}
\textbf{b) (4 P)} Warum ist Zweierkomplement für eine Addierschaltung praktisch?
Erkläre es an einer selbst gewählten Subtraktion mit vier Bits. Zeige die
Bitmuster, die Addition und das gespeicherte Ergebnis.
\antwort{27mm}

\aufgabe{2. Zahlensysteme und Einheiten}{12}
\textbf{a) (8 P)} Ergänze die Tabelle. Alle Werte hier sind \emph{unsigned}.
Verwende genau acht Bits und zwei Hexziffern. Je Lücke 1 Punkt.
\begin{center}
\begin{tabular}{p{35mm}p{55mm}p{35mm}}\toprule
Dezimal&Binär (8 Bits)&Hex\\\midrule
78&\rule{43mm}{0.2pt}&\rule{25mm}{0.2pt}\\[4mm]
\rule{25mm}{0.2pt}&1010 1101&\rule{25mm}{0.2pt}\\[4mm]
\rule{25mm}{0.2pt}&\rule{43mm}{0.2pt}&\code{0xD3}\\[4mm]
\rule{25mm}{0.2pt}&\rule{43mm}{0.2pt}&\code{0xB7}\\\bottomrule
\end{tabular}
\end{center}
\textbf{b) (2 P)} Ein Datenwort hat 24 Bits. Wie viele Bytes und wie viele
Hexziffern sind das?\antwort{8mm}
\textbf{c) (2 P)} Wie viele Bytes sind 1 KiB und 1 kB?\antwort{8mm}

\newpage
\titlepagepart{02 / DARSTELLUNGEN}{Zweierkomplement und Bitbreite}
\aufgabe{3. Zahlenbereiche und Interpretationen}{20}
\textbf{a) (4 P)} Gib jeweils Minimum und Maximum an. Jede Bereichsgrenze zählt
1 Punkt.\par
\begin{center}
\begin{tabular}{lcc}\toprule
Darstellung&Minimum&Maximum\\\midrule
5 Bit unsigned&\rule{35mm}{0.2pt}&\rule{35mm}{0.2pt}\\[5mm]
5 Bit Zweierkomplement&\rule{35mm}{0.2pt}&\rule{35mm}{0.2pt}\\\bottomrule
\end{tabular}
\end{center}
\textbf{b) (6 P)} Werte die Aufrufe aus: Bitmuster in der vorgegebenen Breite,
Grenze $2^{\mathrm{bits}-1}$, ausgeführter Rückgabezweig und Ergebnis.
Je Aufruf 2 Punkte. Rechenweg angeben.
\begin{lstlisting}
from_twos_complement(21, 5)
from_twos_complement(21, 8)
from_twos_complement(128, 8)
\end{lstlisting}
\antwort{34mm}
\textbf{c) (4 P)} Stelle $-22$ mit acht Bits dar. Beginne mit $+22$, invertiere,
addiere 1 und kontrolliere über die signed Stellenwerte. Gib zusätzlich Hex an.
\antwort{28mm}
\textbf{d) (3 P)} Erweitere \code{10101} von fünf auf acht Bits, sodass sein
signed Wert erhalten bleibt. Erkläre deine Ergänzung. Welcher signed Wert
entstünde stattdessen durch Ergänzen von drei Nullen?
\antwort{23mm}
\textbf{e) (3 P)} Warum kann man die kleinste negative Vier-Bit-Zahl nicht als
positive Zahl mit derselben signed Bitbreite darstellen? Zeige, was beim
Invertieren und Addieren von 1 passiert.
\antwort{22mm}

\newpage
\titlepagepart{03 / RECHNEN}{Addition, Subtraktion und Overflow}
\aufgabe{4. Die Hardware rechnet mit acht Bits}{20}
Für \textbf{a--c} gilt: Beide Operanden sind signed. Gib die mathematische
Dezimalsumme, die vollständige Binärsumme, die acht gespeicherten Bits, deren
signed Wert sowie Carry-out und signed Overflow an. Begründe die beiden Flags.
\textbf{Je Fall 5 P:} Dezimalsumme (1), Binärrechnung und gespeicherter Wert (2),
beide Flags mit Begründung (2).

\textbf{a)} \code{0110 1110 + 0010 1101}
\antwort{38mm}
\textbf{b)} \code{1111 1001 + 0000 0111}
\antwort{38mm}
\textbf{c)} \code{1000 1000 + 1110 1100}
\antwort{38mm}
\textbf{d) (5 P)} Berechne $23-41$ durch Addition im Acht-Bit-Zweierkomplement.
Zeige: Negation des zweiten Operanden (2), Binäraddition und gespeicherten
signed Wert (2), signed Overflow mit Begründung (1).
\antwort{35mm}

\newpage
\titlepagepart{04 / PYTHON LESEN}{Schritt für Schritt statt Raten}
\aufgabe{5. Texte, Operatoren und Fehlersuche}{15}
\textbf{a) (5 P)} Was steht nach jeder Zeile in der Variablen? Schreibe
\emph{Textwerte mit Anführungszeichen}, Zahlen ohne. Je Variable 1 Punkt.
\begin{lstlisting}
a = hex(13)
b = a[2:]
c = b.upper()
d = (9 + 3) // 4
e = c.zfill(d)
\end{lstlisting}
\code{a =} \rule{28mm}{0.2pt}\quad\code{b =} \rule{28mm}{0.2pt}\quad
\code{c =} \rule{28mm}{0.2pt}\par
\code{d =} \rule{28mm}{0.2pt}\quad\code{e =} \rule{28mm}{0.2pt}

\textbf{b) (4 P)} Berechne die Ausdrücke. Begründe kurz mit Rechenreihenfolge,
vollständigen Gruppen oder Rest. Je Ausdruck 1 Punkt.
\begin{lstlisting}
(16 + 3) // 4
16 + 3 // 4
19 % 16
-5 % 16
\end{lstlisting}
\antwort{24mm}
\textbf{c) (3 P)} Die Eingabeprüfungen seien korrekt. Finde den Fehler in
folgender Dekodierung. Nenne einen Vier-Bit-Aufruf, bei dem sie falsch rechnet,
das falsche und das richtige Ergebnis sowie die Korrektur.
\begin{lstlisting}
grenze = 2**(bits - 1)
if raw <= grenze:
    return raw
return raw - 2**bits
\end{lstlisting}
\antwort{21mm}
\textbf{d) (3 P)} Erkläre je einen Unterschied:
\begin{enumerate}
\item \code{return ergebnis} und \code{print(ergebnis)};
\item die Zahl \code{9} und der Text \code{"1001"};
\item \code{ValueError} für ungültige Operanden und \code{overflow = True}.
\end{enumerate}
\antwort{20mm}

\newpage
\titlepagepart{05 / SELBST IMPLEMENTIEREN}{Deine vier Funktionen}
\aufgabe{6. Python-Praxisteil}{15}
Erstelle im Tag-4-Ordner eine neue Datei \code{abschluss\_tag04.py}.
Implementiere alle vier Funktionen dort. Du darfst dein Cheatsheet nachschlagen;
schreibe und erkläre den Code selbst. Keine Imports nötig. Ganze Zahlen als
Eingabetypen voraussetzen. Beliebige positive Bitbreiten unterstützen.

\begin{lstlisting}
def to_bin(value: int, bits: int) -> str:
    pass


def to_hex(value: int, bits: int) -> str:
    pass


def from_twos_complement(raw: int, bits: int) -> int:
    pass


def add_with_overflow_flag(a: int, b: int,
                           bits: int) -> tuple[int, bool]:
    pass
\end{lstlisting}
\textbf{Vertrag und Punkte:}
\begin{itemize}
\item \textbf{to\_bin (3 P):} unsigned Eingabe; Text ohne Präfix, genau
\code{bits} Ziffern; links Nullen ergänzen; Bitbreite und Wertebereich prüfen.
\item \textbf{to\_hex (3 P):} unsigned Eingabe; große Hexbuchstaben ohne Präfix;
aufgerundet \code{bits / 4} Stellen; tatsächlichen Bitbereich prüfen, auch wenn
\code{bits} kein Vielfaches von vier ist.
\item \textbf{from\_twos\_complement (4 P):} unsigned Wert eines Musters
entgegennehmen, signed Ganzzahl zurückgeben; gültige Eingaben prüfen;
Vorzeichengrenze korrekt behandeln.
\item \textbf{add\_with\_overflow\_flag (5 P):} signed Operanden prüfen;
richtige Summe bilden; signed Overflow erkennen; Bitbreite begrenzen;
\code{(gespeicherter signed Wert, bool)} zurückgeben.
\end{itemize}
\textbf{Für alle Funktionen:} \code{bits < 1} und Werte außerhalb des jeweiligen
Eingabebereichs lösen \code{ValueError} aus. Ein Overflow der Summe bei gültigen
Operanden löst dagegen keinen Fehler aus.

\block{Erkläre zwei deiner eigenen Codezeilen}
Wähle eine Zeile zur Bitbegrenzung und eine zur Interpretation oder
Overflow-Erkennung. Schreibe jeweils: Was macht sie, und warum brauchen wir sie?
Dies hilft bei der Beurteilung deines Praxiscodes, ohne zusätzliche Punkte.
\antwort{35mm}

\newpage
\titlepagepart{06 / VERIFIKATION}{Tests und persönliche Auswertung}
\aufgabe{7. Beweise, dass dein Code funktioniert}{10}
Erstelle daneben \code{test\_abschluss\_tag04.py}. Importiere die Funktionen aus
\code{abschluss\_tag04}. Schreibe mindestens \textbf{20 ausgeführte Testfälle};
Parametrisierung ist erlaubt. Leite Erwartungen zuerst selbst her.

\textbf{Abdeckung (6 P):} Je 1 Punkt pro Kategorie mit passenden Erwartungen:
\begin{enumerate}
\item Binär/Hex: Null, führende Nullen, oberer unsigned Rand.
\item Hex: eine Bitbreite, die kein Vielfaches von vier ist.
\item Dekodierung: größter positiver Wert, kleinster negativer Wert und $-1$.
\item Addition: ohne Overflow sowie positiver und negativer signed Overflow.
\item Mindestens ein gültiger Ein-Bit-Fall und ein Fall mit mehr als 64 Bits.
\item Ungültige Bitbreite und ungültige unsigned sowie signed Eingaben.
\end{enumerate}
\textbf{Umsetzung (4 P):} Mindestens 20 Fälle (2), korrektes Prüfen von
\code{ValueError} mit \code{pytest.raises} (1), vollständiger Testlauf und
nachvollziehbare Dokumentation (1).

\textbf{Import-Gerüst, keine fertigen Testfälle:}
\begin{lstlisting}
import pytest
from abschluss_tag04 import (
    to_bin, to_hex, from_twos_complement,
    add_with_overflow_flag,
)
\end{lstlisting}
\textbf{Im Terminal aus dem Tag-4-Ordner:}
\begin{lstlisting}[language=bash]
source "$HOME/.venvs/ai-hardware-day04/bin/activate"
python -m pytest -q test_abschluss_tag04.py
\end{lstlisting}
Dieser Aufruf prüft gezielt deine Abschlussdatei. \code{assert} vergleicht
Ist- und Sollwert. Für Fehlerfälle gehört der Funktionsaufruf in den Block
\code{with pytest.raises(ValueError):}.

\textbf{Testlauf:} Datum \rule{30mm}{0.3pt}\quad Bestanden \rule{15mm}{0.3pt}
\quad Fehlgeschlagen \rule{15mm}{0.3pt}

\textbf{Ein Fehler, den ich mit einem Test gefunden habe (oder ein gezielt
abgesicherter Randfall):} Eingabe, Erwartung und Begründung.
\antwort{12mm}

\block{Abgabe und Rückblick}
Gib das ausgefüllte Blatt bzw. deine nummerierten Antworten zusammen mit den
zwei Python-Dateien ab. Markiere, wo du nachgeschlagen hast; das ist kein Fehler.
\par
\textbf{Das Warum kann ich jetzt so erklären:}\antwort{9mm}
\textbf{Hier brauche ich noch eine Erklärung:}\antwort{9mm}
\small Punkte: 1\,/8\quad 2\,/12\quad 3\,/20\quad 4\,/20\quad
5\,/15\quad 6\,/15\quad 7\,/10\qquad Gesamt: \rule{12mm}{0.3pt}/100.
\par Grundlage: Lernplan Woche 1, Tag 4; DDCA RISC-V Edition, Abschnitt 1.4;
unsere Python-Übung. Die Punkte dienen der Rückmeldung, nicht einer Schulnote.
\end{document}
